%=====================================================================
% INSTALACOES ELETRICAS — SEIS PROJETOS COMPLETOS
% 2 residenciais | 2 comerciais | 2 industriais
% ABNT + Equatorial Maranhao
%=====================================================================
\documentclass[11pt,aspectratio=169]{beamer}
\usepackage[utf8]{inputenc}
\usepackage[T1]{fontenc}
\usepackage[brazil]{babel}
\usepackage{lmodern}
\usepackage{amsmath,amssymb,bm}
\usepackage{siunitx}
\usepackage{booktabs,array,tabularx,multirow}
\usepackage[most]{tcolorbox}
\usepackage{tikz}
\usepackage{pgfplots}
\usepackage[european,american currents]{circuitikz}
\usepackage{ragged2e}
\usepackage{microtype}
\usepackage{xcolor}
\usetikzlibrary{arrows.meta,positioning,calc,fit,shapes.geometric,patterns,decorations.pathmorphing,decorations.markings}
\pgfplotsset{compat=1.18}
\sisetup{locale=DE,per-mode=symbol,detect-all}
\newcommand{\Disciplina}{Instalações Elétricas Prediais e Industriais}
\newcommand{\TituloCurto}{6 Projetos}
%=====================================================================
% ESTILO DA CASA -- TEMPLATE BEAMER (padrao da colecao)
% Arquivo central de identidade visual para todos os materiais.
% Azul-petroleo + ambar | tcolorbox | TikZ/pgfplots/circuitikz
%=====================================================================

% Dependencias minimas do proprio estilo.
\RequirePackage{xcolor}
\RequirePackage{tikz}
\usetikzlibrary{calc}
\RequirePackage{tcolorbox}
\tcbuselibrary{skins,breakable}
\RequirePackage{listings}
\RequirePackage{xparse}
\RequirePackage{etoolbox}

\makeatletter

%====================== PALETA OFICIAL ================================
\definecolor{corPrim}{RGB}{13,48,82}      % azul-petroleo
\definecolor{corAccent}{RGB}{222,138,28}  % ambar
\definecolor{corDef}{RGB}{31,111,168}     % azul medio
\definecolor{corNorma}{RGB}{0,121,107}    % verde-petroleo
\definecolor{corEx}{RGB}{46,125,50}       % verde
\definecolor{corFix}{RGB}{204,102,0}      % laranja
\definecolor{corAt}{RGB}{176,58,46}       % vermelho
\definecolor{codBg}{RGB}{252,250,245}
\definecolor{codKw}{RGB}{13,48,82}
\definecolor{codStr}{RGB}{46,125,50}
\definecolor{codCom}{RGB}{120,130,138}

% Compatibilidade com nomes usados historicamente nos materiais.
\colorlet{Petroleo}{corPrim}
\colorlet{PetroleoEscuro}{corPrim!82!black}
\colorlet{PetroleoMedio}{corDef!82!corPrim}
\colorlet{PetroleoClaro}{corDef!7}
\colorlet{Ambar}{corAccent}
\colorlet{AmbarEscuro}{corAccent!72!black}
\colorlet{AmbarClaro}{corAccent!16}
\colorlet{AzulPetroleo}{corPrim}
\colorlet{AzulPetroleoEscuro}{corPrim!82!black}
\colorlet{AzulPetroleoClaro}{corDef!7}
\colorlet{AzulMedio}{corDef}
\colorlet{AzulClaro}{corDef!7}
\colorlet{CinzaEscuro}{black!82}
\colorlet{CinzaClaro}{black!5}
\colorlet{CinzaSuave}{black!3}
\colorlet{CinzaMedio}{black!55}
\colorlet{CinzaTexto}{black!82}
\colorlet{CinzaBox}{black!3}
\colorlet{Verde}{corEx}
\colorlet{Vermelho}{corAt}
\colorlet{Turquesa}{corNorma!70!corDef}
\colorlet{Roxo}{corDef!58!corAt}
\colorlet{ifmagreen}{corPrim}
\colorlet{ifmagreen2}{corDef}
\colorlet{ifmalight}{corDef!7}
\colorlet{accent}{corAccent}
\colorlet{warn}{corAt}
\colorlet{ok}{corEx}
\colorlet{bluex}{corDef}
\colorlet{grayx}{black!58}

% Variantes em minusculas presentes em materiais antigos.
\colorlet{azulPetroleo}{corPrim}
\colorlet{azulPetroleoEscuro}{corPrim!82!black}
\colorlet{azulPetroleoClaro}{corDef!7}
\colorlet{azulMedio}{corDef}
\colorlet{azulClaro}{corDef!7}
\colorlet{azulEscuro}{corPrim!82!black}
\colorlet{azul}{corPrim}
\colorlet{azul2}{corDef}
\colorlet{ambar}{corAccent}
\colorlet{ambarClaro}{corAccent!16}
\colorlet{cinza}{black!5}
\colorlet{cinzaClaro}{black!5}
\colorlet{cinzaMedio}{black!55}
\colorlet{cinzaEscuro}{black!82}
\colorlet{cinzaTexto}{black!82}
\colorlet{verde}{corEx}
\colorlet{verdeclaro}{corEx!10}
\colorlet{vermelho}{corAt}
\colorlet{vermelhoclaro}{corAt!9}
\colorlet{laranja}{corFix}
\colorlet{roxo}{corDef!58!corAt}

%====================== TEMA BEAMER ==================================
\usetheme{default}
\usefonttheme{professionalfonts}
\setbeamercolor{normal text}{fg=black!85,bg=white}
\setbeamercolor{structure}{fg=corPrim}
\setbeamercolor{frametitle}{bg=corPrim,fg=white}
\setbeamercolor{footsec}{bg=corPrim,fg=white!92!corPrim}
\setbeamercolor{alerted text}{fg=corAccent!85!black}
\setbeamerfont{frametitle}{series=\bfseries,size=\large}
\setbeamerfont{footline}{size=\tiny}
\setbeamertemplate{navigation symbols}{}
\setbeamertemplate{caption}{\raggedright\insertcaption\par}
\setbeamersize{text margin left=0.55cm,text margin right=0.55cm}

\setbeamertemplate{itemize item}{\raise0.35ex\hbox{\color{corAccent}\rule{0.62ex}{0.62ex}}}
\setbeamertemplate{itemize subitem}{\color{corDef}\raise0.5pt\hbox{--}}
\setbeamertemplate{itemize subsubitem}{\color{corPrim}\tiny$\circ$}
\setbeamertemplate{enumerate item}{\color{corPrim}\textbf{\insertenumlabel.}}
\setbeamertemplate{section in toc}{%
  \leavevmode\textcolor{corAccent}{\textbf{\inserttocsectionnumber.}}~\inserttocsection\par\vskip1.5pt}
\setbeamertemplate{subsection in toc}{%
  \leavevmode\leftskip=0.95cm\textcolor{corDef}{\textbullet}~\inserttocsubsection\par}

%====================== TITULO DO SLIDE ==============================
\setbeamertemplate{frametitle}{%
  \nointerlineskip%
  \begin{beamercolorbox}[wd=\paperwidth,ht=2.7ex,dp=1.3ex,leftskip=0.55cm,rightskip=0.55cm]{frametitle}%
    \usebeamerfont{frametitle}\insertframetitle%
  \end{beamercolorbox}%
  {\color{corAccent}\rule{\paperwidth}{1.3pt}}\par\vskip2pt%
}

%====================== RODAPE DUPLO ================================
\newcommand{\CasaRodapeEsq}{%
  \ifcsname Disciplina\endcsname\Disciplina\else\insertshorttitle\fi}
\newcommand{\CasaRodapeDir}{%
  \ifcsname TituloCurto\endcsname\TituloCurto\else\insertshorttitle\fi}
\newcommand{\CasaMetaCapa}{%
  \ifcsname Disciplina\endcsname
    \if\relax\detokenize\expandafter{\Disciplina}\relax
    \else
      \Disciplina\ \textbullet\ \insertdate
    \fi
  \else
    \insertdate
  \fi}
\setbeamertemplate{footline}{%
  \leavevmode%
  \hbox{%
    \begin{beamercolorbox}[wd=.5\paperwidth,ht=2.4ex,dp=1ex,leftskip=0.5cm]{footsec}%
      \usebeamerfont{footline}\CasaRodapeEsq%
    \end{beamercolorbox}%
    \begin{beamercolorbox}[wd=.5\paperwidth,ht=2.4ex,dp=1ex,rightskip=0.5cm]{footsec}%
      \hfill\usebeamerfont{footline}\CasaRodapeDir\quad\textbar\quad\insertframenumber%
    \end{beamercolorbox}%
  }\vskip0pt%
}

%====================== DIVISORIAS DE SECAO =========================
\newcommand{\CasaSecNum}{\ifnum\value{section}<10 0\fi\arabic{section}}
\newcommand{\CasaDivisoriaManual}[2]{%
  \begin{frame}[plain,noframenumbering]
    \makebox[0pt][l]{%
      \raisebox{0pt}[0pt][0pt]{%
        \begin{tikzpicture}[remember picture,overlay]
          \fill[corPrim] (current page.south west) rectangle (current page.north east);
          \if\relax\detokenize{#1}\relax\else
            \node[anchor=west,text=corAccent,font=\bfseries\fontsize{40}{40}\selectfont]
              at ([xshift=1.6cm,yshift=1.1cm]current page.west) {#1};
          \fi
          \fill[corAccent] ([xshift=1.7cm,yshift=0.35cm]current page.west)
            rectangle ++(2.4,0.09);
          \node[anchor=west,text=white,font=\bfseries\LARGE,text width=0.75\paperwidth]
            at ([xshift=1.6cm,yshift=-0.5cm]current page.west) {#2};
        \end{tikzpicture}%
      }%
    }%
  \end{frame}%
}
\AtBeginSection[]{\CasaDivisoriaManual{\CasaSecNum}{\insertsection}}

% Alguns materiais antigos usam \secslide{titulo}; outros, \secslide{numero}{titulo}.
\@ifundefined{secslide}
  {\NewDocumentCommand{\secslide}{m g}{\IfNoValueTF{#2}{\CasaDivisoriaManual{}{#1}}{\CasaDivisoriaManual{#1}{#2}}}}
  {\RenewDocumentCommand{\secslide}{m g}{\IfNoValueTF{#2}{\CasaDivisoriaManual{}{#1}}{\CasaDivisoriaManual{#1}{#2}}}}
\@ifundefined{Divisoria}
  {\NewDocumentCommand{\Divisoria}{m g}{\IfNoValueTF{#2}{\CasaDivisoriaManual{}{#1}}{\CasaDivisoriaManual{#1}{#2}}}}
  {\RenewDocumentCommand{\Divisoria}{m g}{\IfNoValueTF{#2}{\CasaDivisoriaManual{}{#1}}{\CasaDivisoriaManual{#1}{#2}}}}

%====================== CAPA ========================================
\defbeamertemplate*{title page}{estilo da casa}{%
  \vspace{0.8cm}%
  {\color{white}\bfseries\huge\inserttitle\par}%
  \vspace{2mm}%
  {\color{corAccent}\rule{3.2cm}{2.2pt}\par}%
  \vspace{4mm}%
  {\color{white!85}\large\insertsubtitle\par}%
  \vspace{8mm}%
  {\color{white}\normalsize\insertauthor\par}%
  \vspace{1mm}%
  {\color{white!70}\small\CasaMetaCapa\par}%
}
\setbeamertemplate{title page}[estilo da casa]

%====================== CAIXAS =======================================
% Materiais antigos usam convencoes diferentes para o argumento opcional
% das caixas. Para preservar a semantica, caixas existentes nao sao
% redefinidas; apenas sua paleta historica e recolorida acima.
\tcbset{
  caixabase/.style={
    enhanced,sharp corners=downhill,boxrule=0.6pt,left=2mm,right=2mm,
    top=1.2mm,bottom=1.2mm,fonttitle=\bfseries\small,
    attach boxed title to top left={xshift=3mm,yshift=-2.4mm},
    boxed title style={sharp corners,boxrule=0pt},
    before skip=3pt,after skip=3pt
  }
}
\@ifundefined{caixaDef}{\newtcolorbox{caixaDef}[1][]{caixabase,colback=corDef!5,colframe=corDef,coltitle=white,colbacktitle=corDef,title=Definição,#1}}{}
\@ifundefined{caixaNorma}{\newtcolorbox{caixaNorma}[1][]{caixabase,colback=corNorma!5,colframe=corNorma,coltitle=white,colbacktitle=corNorma,title=Norma / Método,#1}}{}
\@ifundefined{caixaEx}{\newtcolorbox{caixaEx}[1][]{caixabase,colback=corEx!5,colframe=corEx,coltitle=white,colbacktitle=corEx,title=Exemplo resolvido,#1}}{}
\@ifundefined{caixaFix}{\newtcolorbox{caixaFix}[1][]{caixabase,colback=corFix!5,colframe=corFix,coltitle=white,colbacktitle=corFix,title=Fixação,#1}}{}
\@ifundefined{caixaAt}{\newtcolorbox{caixaAt}[1][]{caixabase,colback=corAt!5,colframe=corAt,coltitle=white,colbacktitle=corAt,title=Atenção,#1}}{}
\@ifundefined{caixaForm}{\newtcolorbox{caixaForm}[1][]{caixabase,colback=corPrim!4,colframe=corPrim,coltitle=white,colbacktitle=corPrim,title=Formulário,#1}}{}
\@ifundefined{caixaObs}{\newtcolorbox{caixaObs}[1][]{caixabase,colback=black!3,colframe=corPrim,coltitle=white,colbacktitle=corPrim,title=Observação,#1}}{}
\@ifundefined{caixaAlerta}{\newtcolorbox{caixaAlerta}[1][]{caixabase,colback=corAt!5,colframe=corAt,coltitle=white,colbacktitle=corAt,title=Atenção,#1}}{}
\@ifundefined{caixaProjeto}{\newtcolorbox{caixaProjeto}[1][]{caixabase,colback=corNorma!5,colframe=corNorma,coltitle=white,colbacktitle=corNorma,title=Projeto prático,#1}}{}
\@ifundefined{caixaProj}{\newtcolorbox{caixaProj}[1][]{caixabase,colback=corNorma!5,colframe=corNorma,coltitle=white,colbacktitle=corNorma,title=Projeto prático,#1}}{}

%====================== CODIGO =======================================
\lstdefinestyle{codCasa}{
  basicstyle=\ttfamily\scriptsize,
  keywordstyle=\color{codKw}\bfseries,
  stringstyle=\color{codStr},commentstyle=\color{codCom}\itshape,
  numbers=left,numberstyle=\tiny\color{codCom},numbersep=5pt,
  frame=single,rulecolor=\color{corPrim!25},backgroundcolor=\color{codBg},
  breaklines=true,showstringspaces=false,tabsize=2,columns=fullflexible
}
\lstdefinestyle{codC}{style=codCasa,language=C}
\lstdefinestyle{codPy}{style=codCasa,language=Python}
\lstdefinestyle{py}{style=codCasa,language=Python}

%====================== EXTENSOES MODULARES ==========================
% Se existir <jobname>_antes_laboratorios.tex, o arquivo e carregado uma
% unica vez imediatamente antes da secao de laboratorios. Isso permite
% ampliar um material sem duplicar o arquivo principal.
\newif\ifCasaExtensaoAntesLabsInserida
\CasaExtensaoAntesLabsInseridafalse
\AtBeginDocument{%
  \IfFileExists{\jobname_antes_laboratorios.tex}{%
    \let\CasaSectionOriginal\section
    \renewcommand{\section}[1]{%
      \ifCasaExtensaoAntesLabsInserida\else
        \ifstrequal{##1}{Laboratórios, projetos e comissionamento}{%
          \global\CasaExtensaoAntesLabsInseridatrue
          \input{\jobname_antes_laboratorios.tex}%
        }{}%
      \fi
      \CasaSectionOriginal{##1}%
    }%
  }{}%
}

\makeatother

\title{Instalações Elétricas — Seis Projetos Completos}
\subtitle{2 residenciais, 2 comerciais e 2 industriais — ABNT + Equatorial Maranhão}
\author{Coleção de Engenharia Elétrica}
\institute{Material 03 — caderno de projetos}
\date{Atualização: agosto de 2026}
\begin{document}
{\setbeamercolor{background canvas}{bg=corPrim}\setbeamertemplate{footline}{}\begin{frame}[plain,noframenumbering]\titlepage\end{frame}}
\begin{frame}{Objetivo deste caderno}
\small
\begin{columns}[T,onlytextwidth]
\column{0.49\textwidth}
\begin{caixaProjeto}[title=Seis arquiteturas diferentes]
\begin{itemize}
\item casa térrea com duas suítes;
\item apartamento com varanda e duas suítes;
\item edifício corporativo com elevador, escada e SPDA;
\item clínica/centro de serviços;
\item indústria leve atendida em BT;
\item indústria com subestação própria e SPDA.
\end{itemize}
\end{caixaProjeto}
\column{0.49\textwidth}
\begin{caixaNorma}[title=Base técnica]
Os casos usam as normas ABNT aplicáveis e os documentos vigentes da Equatorial para BT, múltiplas unidades, média tensão e subestações. As plantas são bases didáticas funcionais para o projeto elétrico; obra real exige projeto arquitetônico, levantamentos, dados de placa, análise de risco, dados da rede e responsabilidade técnica.
\end{caixaNorma}
\end{columns}
\end{frame}
% Agregador canônico — seis projetos com plantas arquitetônicas funcionais
% 2 residenciais | 2 comerciais | 2 industriais
\input{projetos/base/simbologia_eletrica.tex}
\input{projetos/base/plantas_arquitetonicas.tex}
% Portas e vaos complementares para as plantas realistas
% Implementacao simples e robusta, sem aritmetica adicional nas coordenadas TikZ.
\newcommand{\ArqVaoH}[3]{%
  \draw[white,line width=5.2pt] (#1,#3)--(#2,#3);
  \draw[black!20,thin] (#1,#3)--(#2,#3);}
\newcommand{\ArqVaoV}[3]{%
  \draw[white,line width=5.2pt] (#3,#1)--(#3,#2);
  \draw[black!20,thin] (#3,#1)--(#3,#2);}
\newcommand{\ArqPortaCorrerH}[3]{%
  \draw[white,line width=5.4pt] (#1,#3)--(#2,#3);
  \draw[black!55,thin,double=white,double distance=1.35pt] (#1,#3)--(#2,#3);}
\newcommand{\ArqPortaCorrerV}[3]{%
  \draw[white,line width=5.4pt] (#3,#1)--(#3,#2);
  \draw[black!55,thin,double=white,double distance=1.35pt] (#3,#1)--(#3,#2);}
\newcommand{\ArqPortaElevV}[3]{%
  \draw[white,line width=5.0pt] (#3,#1)--(#3,#2);
  \draw[black!65,thin,double=white,double distance=1.15pt] (#3,#1)--(#3,#2);}


\input{projetos/residenciais/casa_terrea.tex}
% Residencial II - apartamento com varanda, duas suites e acessos coerentes
\section{Projeto Residencial II — apartamento com varanda}

\newcommand{\AptoFinalBase}{%
\draw[arqExt] (0,0) rectangle (12,9);\draw[arqExt] (12,4) rectangle (14,9);
\ArqParedeInt{4.2}{0}{4.2}{9}\ArqParedeInt{5.2}{2}{5.2}{6.2}\ArqParedeInt{0}{4.4}{4.2}{4.4}
\ArqParedeInt{2.6}{0}{2.6}{2.1}\ArqParedeInt{2.6}{2.1}{4.2}{2.1}
\ArqParedeInt{2.6}{6.8}{4.2}{6.8}\ArqParedeInt{2.6}{6.8}{2.6}{9}
\ArqParedeInt{5.2}{3.8}{12}{3.8}\ArqParedeInt{7.2}{0}{7.2}{3.8}\ArqParedeInt{10}{0}{10}{3.8}
\ArqParedeInt{5.2}{2.2}{6.6}{2.2}\ArqParedeInt{6.6}{2.2}{6.6}{3.8}
% entrada e acessos
\ArqPortaHUp{5.7}{0}{1.0}\node[arqSub] at (6.2,0.45){ENTRADA / HALL COMUM};
\ArqVaoH{5.65}{7.05}{3.8} % foyer -> sala
\ArqVaoH{7.65}{9.45}{3.8} % cozinha -> sala
\ArqPortaVLeft{5.2}{2.45}{0.85} % foyer -> hall intimo
\ArqPortaVLeft{4.2}{2.55}{0.85}\ArqPortaVLeft{4.2}{4.90}{0.85} % suites
% banhos compactos: portas de correr evitam conflito com loucas e box
\ArqPortaCorrerH{2.95}{3.75}{2.1}\ArqPortaCorrerH{2.95}{3.75}{6.8}
\ArqPortaHUp{5.45}{2.2}{0.75} % lavabo abre para o foyer
\ArqPortaVRight{10}{1.10}{0.80} % cozinha -> servico
\ArqPortaCorrerV{5.0}{8.2}{12} % sala -> varanda
% janelas
\ArqJanV{0.55}{3.4}{0}\ArqJanV{5.25}{8.35}{0}\ArqJanH{7.65}{9.4}{0}\ArqJanH{10.35}{11.6}{0}\ArqJanH{0.6}{2.3}{9}
% suite 2
\ArqCama{0.45}{0.50}{1.95}{2.55}\ArqArmario{0.4}{3.45}{2.8}{0.5}\ArqChuveiro{2.78}{0.15}{1.15}{0.78}\ArqVaso{3.55}{1.35}\ArqPia{2.72}{1.15}{0.55}{0.4}\node[arqLabel] at (1.55,3.9){SUÍTE 2};\node[arqSub] at (3.45,1.92){B2};
% suite master
\ArqCama{0.45}{4.95}{2.05}{2.70}\ArqArmario{0.45}{8.05}{1.85}{0.5}\ArqChuveiro{2.78}{8.05}{1.15}{0.78}\ArqVaso{3.55}{7.35}\ArqPia{2.72}{7}{0.55}{0.4}\node[arqLabel] at (1.5,8.55){SUÍTE MASTER};\node[arqSub] at (3.45,8.72){B1};
\node[arqSub,rotate=90] at (4.70,4.25){HALL ÍNTIMO};
% foyer e lavabo
\node[arqSub] at (6.2,1.45){FOYER};\ArqVaso{5.75}{2.85}\ArqPia{5.35}{3.25}{0.55}{0.38}\node[arqSub] at (5.9,3.6){LAVABO};
% cozinha e servico
\ArqBancada{7.45}{0.25}{2.15}{0.55}\ArqBancada{9.35}{0.8}{0.55}{2.5}\ArqFogao{8.55}{0.18}\ArqGeladeira{7.35}{2.7}\node[arqLabel] at (8.55,3.35){COZINHA};
\ArqLavadora{10.35}{0.35}\ArqBancada{11.1}{0.3}{0.55}{2.65}\node[arqLabel] at (11.05,3.35){SERVIÇO};
% estar e varanda
\ArqSofa{5.8}{4.55}{3.2}{1.05}\ArqBancada{5.55}{7.75}{0.38}{1.0}\ArqMesa{8.65}{6.35}{2.15}{1.0}\node[arqLabel] at (8.65,5.95){ESTAR / JANTAR};
\ArqSofa{12.45}{5.05}{1.1}{2.2}\ArqMesa{12.55}{7.65}{0.85}{0.7}\node[arqLabel,rotate=90] at (13.55,6.5){VARANDA};
\CotaH{0}{12}{9.18}{12,00 m}\CotaV{0}{9}{14.18}{9,00 m}
}

\begin{frame}{Residencial II — acessos refinados}
\centering\begin{tikzpicture}[scale=0.53,every node/.style={font=\tiny}]\AptoFinalBase\end{tikzpicture}
\vspace{0.4mm}\scriptsize Entrada pelo hall comum; lavabo voltado ao foyer; cozinha integrada à sala; serviço acessado pela cozinha; B1/B2 por dentro das suítes, com portas de correr; varanda pela sala.
\end{frame}

\begin{frame}{Residencial II — iluminação A1/A2}
\centering\begin{tikzpicture}[scale=0.53,every node/.style={font=\tiny}]
\AptoFinalBase\SimQD{6.85}{3.30}{QD-A}
\foreach \x/\y in {1.45/2.8,3.35/1,1.5/6,3.35/8,4.7/4.3,6.2/1.2,5.9/2.9,8.5/1.9,11/1.9,8/5.1,10.1/6.5,13/6.6}{\SimLuz{\x}{\y}}
\foreach \x/\y in {3.92/2.82,3.12/2.32,3.92/5.12,3.12/6.58,5.42/0.42,5.42/2.55,6.82/3.52,9.82/1.28,5.42/4.18,11.72/5.18}{\SimInterruptor{\x}{\y}{S}}
\draw[corAccent!80!black,dashed,thick] (6.85,3.30)--(8.0,5.1)--(10.1,6.5)--(13,6.6);
\draw[corAccent!80!black,dashed,thick] (6.85,3.30)--(8.5,1.9)--(11,1.9);
\draw[corAccent!55!black,dashed,thick] (6.85,3.30)--(6.2,1.2)--(5.9,2.9)--(4.7,4.3)--(1.5,6)--(3.35,8);
\draw[corAccent!55!black,dashed,thick] (4.7,4.3)--(1.45,2.8)--(3.35,1);
\node[font=\tiny\bfseries,fill=white,inner sep=1pt,corAccent!80!black] at (9.5,4.8){A1};
\node[font=\tiny\bfseries,fill=white,inner sep=1pt,corAccent!55!black] at (3.2,4.3){A2};
\end{tikzpicture}
\vspace{0.3mm}\scriptsize Previsão: A1 = 0,94 kVA e A2 = 0,94 kVA. Os comandos estão junto aos acessos; a potência mínima total é 1,88 kVA, calculada pelas áreas e pelos acréscimos completos de 4 m$^2$.
\end{frame}

\begin{frame}{Residencial II — 27 TUG e sete cargas específicas}
\centering\begin{tikzpicture}[scale=0.53,every node/.style={font=\tiny}]
\AptoFinalBase\SimQD{6.85}{3.30}{QD-A}
% A3: suites, hall, foyer, sala e varanda — 16 pontos
\foreach \x/\y in {0.25/0.55,0.25/4.15,2.35/0.25,3.95/4.15,0.25/4.85,0.25/8.75,2.35/8.75,3.95/4.65,4.35/4.30,5.35/0.55,5.45/4.05,7.40/4.05,9.50/4.05,11.70/4.10,11.70/8.70,13.70/8.70}{\SimTUG{\x}{\y}}
% A4 cozinha, A5 servico, A6 banheiros
\foreach \x/\y in {7.45/0.25,8.55/0.25,9.75/1.50,7.35/3.55}{\SimTUGM{\x}{\y}}
\foreach \x/\y in {10.15/0.25,11.75/0.25,10.15/3.55,11.75/3.55}{\SimTUGM{\x}{\y}}
\foreach \x/\y in {3.95/1.65,3.95/7.15,5.45/3.50}{\SimTUG{\x}{\y}}
% A7--A13
\SimTUE{3.4}{0.4}{A7}\SimTUE{3.4}{8.55}{A8}\SimTUE{8.9}{0.35}{A9}\SimTUE{10.7}{0.45}{A10}
\SimTUE{0.3}{3.1}{A11}\SimTUE{0.3}{7.95}{A12}\SimTUE{11.55}{7.8}{A13}
\end{tikzpicture}
\vspace{0.3mm}\tiny A3 áreas secas/varanda (16) · A4 cozinha (4) · A5 serviço (4) · A6 banheiros/lavabo (3) · A7--A8 chuveiros · A9 forno · A10 lavadora · A11--A13 ar-condicionado.
\end{frame}

\begin{frame}[shrink=8]{Residencial II — memória dos mínimos por ambiente}
\scriptsize
\begin{tabularx}{\textwidth}{>{\bfseries}p{2.35cm}c c c c >{\centering\arraybackslash}X}
\toprule
Ambiente & Área (m$^2$) & Perím. (m) & Luz (VA) & TUG & TUG (VA)\\
\midrule
Suíte 2 / B2 & 15,1 / 3,4 & 17,2 / próprio & 220 / 100 & 4 / 1 & 400 / 600\\
Suíte master / B1 & 15,8 / 3,5 & 17,6 / próprio & 220 / 100 & 4 / 1 & 400 / 600\\
Hall íntimo & 4,2 & critério por área & 100 & 1 & 100\\
Foyer & 5,4 & critério por área & 100 & 1 & 100\\
Lavabo & 2,2 & critério próprio & 100 & 1 & 600\\
Sala/jantar & 35,4 & 24,0 & 520 & 5 & 500\\
Cozinha & 10,6 & 13,2 & 160 & 4 & 1.900\\
Serviço & 7,6 & 11,6 & 100 & 4 & 1.900\\
Varanda & 10,0 & ao menos uma & 160 & 1 & 100\\
\midrule
\textbf{Total} & -- & -- & \textbf{1.880} & \textbf{27} & \textbf{7.200}\\
\bottomrule
\end{tabularx}
\end{frame}

\begin{frame}[shrink=6]{Residencial II — circuitos A1 a A6}
\scriptsize
\begin{tabularx}{\textwidth}{c p{3.5cm} c c c c c X}
\toprule
C & Uso & $S$ & $I_b$ & Cu & DJ & Fase & Solução adotada\\
\midrule
A1 & Luz social/cozinha/serviço/varanda & 0,94 kVA & 4,3 A & 1,5 mm$^2$ & 10 A & A & RCBO 30 mA\\
A2 & Luz íntima/foyer/lavabo & 0,94 kVA & 4,3 A & 1,5 mm$^2$ & 10 A & B & RCBO 30 mA\\
A3 & TUG áreas secas e varanda & 1,60 kVA & 7,3 A & 2,5 mm$^2$ & 16 A & C & RCBO 30 mA\\
A4 & TUG cozinha & 1,90 kVA & 8,6 A & 2,5 mm$^2$ & 16 A & B & RCBO 30 mA\\
A5 & TUG serviço & 1,90 kVA & 8,6 A & 2,5 mm$^2$ & 16 A & C & RCBO 30 mA\\
A6 & TUG banheiros/lavabo & 1,80 kVA & 8,2 A & 2,5 mm$^2$ & 16 A & C & RCBO 30 mA\\
\bottomrule
\end{tabularx}
\end{frame}

\begin{frame}[shrink=6]{Residencial II — circuitos A7 a A13}
\scriptsize
\begin{tabularx}{\textwidth}{c p{3.1cm} c c c c c X}
\toprule
C & Equipamento & $P$ & $I_b$ & Cu & DJ & Fase & Proteção/observação\\
\midrule
A7 & Chuveiro B2 & 4,50 kW & 20,5 A & 4 mm$^2$ & 25 A & A & dedicado; RCBO 30 mA\\
A8 & Chuveiro B1 & 4,50 kW & 20,5 A & 4 mm$^2$ & 25 A & B & dedicado; RCBO 30 mA\\
A9 & Forno elétrico & 3,50 kW & 15,9 A & 4 mm$^2$ & 20 A & C & dedicado; RCBO 30 mA\\
A10 & Lavadora & 1,20 kW & 5,5 A & 2,5 mm$^2$ & 16 A & B & dedicado; RCBO 30 mA\\
A11--A13 & Ar-condicionado & 3$\times$1,00 kW & 4,5 A cada & 2,5 mm$^2$ & 16 A & A/A/A & dedicados; fabricante\\
\bottomrule
\end{tabularx}
\begin{caixaAt}
Seleções iniciais: recalcular $I_z$, agrupamento, queda, curto, capacidade de interrupção e requisitos do fabricante. A proteção diferencial adotada deve ser compatível com cargas eletrônicas/inverter.
\end{caixaAt}
\end{frame}

\begin{frame}{Residencial II — balanceamento e interface EMUC}
\small\begin{columns}[T,onlytextwidth]
\column{0.49\textwidth}
\begin{caixaProjeto}[title=Unidade revisada]
\begin{align*}
S_A&=\SI{8.44}{kVA}\;(\SI{38.4}{A})\\
S_B&=\SI{8.54}{kVA}\;(\SI{38.8}{A})\\
S_C&=\SI{8.80}{kVA}\;(\SI{40.0}{A})
\end{align*}
Total equivalente: \SI{25.78}{kW}. Proteção de 63 A é a premissa inicial da unidade. Para esta ilustração, os VA de previsão foram tomados numericamente como W; a declaração real usa potências ativas nominais. Verificar também cenários de demanda, pois o balanceamento instalado concentra os três aparelhos de ar-condicionado na fase A.
\end{caixaProjeto}
\column{0.49\textwidth}
\begin{caixaNorma}[title=Edifício multifamiliar]
Centro de medição, demanda do conjunto, prumadas, queda, proteção geral e solução de conexão são definidos pela \textbf{NT.00004.EQTL Rev.08} e seu Anexo I vigente. Não multiplicar a carga desta unidade pelo número de apartamentos sem aplicar o método do empreendimento.
\end{caixaNorma}
\end{columns}
\end{frame}

\begin{frame}{Residencial II — unifilar da unidade}
\centering\begin{tikzpicture}[scale=0.74,every node/.style={font=\tiny}]
\tikzset{b/.style={draw=corPrim,fill=corDef!8,rounded corners,minimum width=2cm,minimum height=0.56cm,align=center}}
\node[b](cm)at(0,4.4){Centro de medição\\EMUC};\node[b](pr)at(0,3.5){Prumada 3F+N+PE};\node[b](qd)at(0,2.55){QD-A\\DG 3P 63 A + DPS};
\node[b](l)at(-3.6,0.9){A1--A6\\luz/TUG + RCBO};\node[b](ch)at(-1.2,0.9){A7--A8\\chuveiros};\node[b](co)at(1.2,0.9){A9--A10\\forno/lavadora};\node[b](ac)at(3.6,0.9){A11--A13\\climatização};
\node[b,draw=corNorma,fill=corNorma!7](bep)at(4.8,2.9){BEP do edifício\\PE/equipot.};
\draw[-{Stealth},thick,corPrim](cm)--(pr)--(qd);\foreach \x in {l,ch,co,ac}{\draw[-{Stealth},thick](qd)--(\x);}\draw[corNorma,thick](pr)--(bep);\draw[corNorma,thick](qd)--(bep);\draw[corNorma,thick](bep.south)--(4.8,1.55)node[ground]{};
\end{tikzpicture}\end{frame}

% Comercial I - edificio corporativo com nucleo, acessos e SPDA
\section{Projeto Comercial I — edifício corporativo com SPDA}

\newcommand{\EdifComBase}{%
\draw[arqExt] (0,0) rectangle (26,18);
\ArqParedeInt{0}{8}{26}{8}\ArqParedeInt{0}{10}{26}{10}
\foreach \x in {5.2,10.4,15.6,20.8}{\ArqParedeInt{\x}{10}{\x}{18}}
\ArqParedeInt{5}{0}{5}{8}\ArqParedeInt{10}{0}{10}{8}\ArqParedeInt{18}{0}{18}{8}\ArqParedeInt{22}{0}{22}{8}
% nucleo: escada, elevador, shaft e sanitarios com corredor lateral de 1,30 m
\ArqParedeInt{13.3}{0}{13.3}{8}\ArqParedeInt{16.7}{0}{16.7}{8}
\ArqParedeInt{14.3}{0}{14.3}{5.8}
\ArqParedeInt{14.3}{1.8}{16.7}{1.8}\ArqParedeInt{14.3}{3.7}{16.7}{3.7}\ArqParedeInt{14.3}{5.8}{16.7}{5.8}
% portas salas superiores e inferiores
\foreach \x in {1.0,6.2,11.4,16.6,21.8}{\ArqPortaHUp{\x}{10}{0.90}}
\foreach \x in {1.0,6.0,18.7,22.7}{\ArqPortaHDown{\x}{8}{0.90}}
% escada, elevador e shaft acessados diretamente pelo corredor central
\ArqPortaHDown{11.2}{8}{0.90}\ArqPortaCorrerH{13.75}{15.25}{8}\ArqPortaHDown{15.75}{8}{0.70}
% corredor dos sanitarios: acesso independente do patamar da escada
\ArqVaoH{16.85}{17.95}{8}
\ArqPortaVLeft{16.7}{0.40}{0.75}\ArqPortaVLeft{16.7}{2.25}{0.75}\ArqPortaVLeft{16.7}{4.25}{0.90}
% fachada
\foreach \xa/\xb in {0.8/4.3,5.9/9.5,11.1/14.7,16.3/19.9,21.5/25.2}{\ArqJanH{\xa}{\xb}{18}}
\foreach \xa/\xb in {0.8/4.2,5.8/9.2,18.5/21.4,22.5/25.2}{\ArqJanH{\xa}{\xb}{0}}
% mobiliario salas superiores
\foreach \x in {0.6,5.8,11.0,16.2,21.4}{\ArqDesk{\x}{12.0}{1.45}{0.60}\ArqDesk{\x+2.3}{12.0}{1.45}{0.60}\ArqDesk{\x}{15.2}{1.45}{0.60}\ArqDesk{\x+2.3}{15.2}{1.45}{0.60}}
% salas inferiores
\foreach \x in {0.5,5.5,18.4}{\ArqDesk{\x}{2.0}{1.25}{0.58}\ArqDesk{\x+2.0}{2.0}{1.25}{0.58}\ArqDesk{\x}{5.2}{1.25}{0.58}\ArqDesk{\x+2.0}{5.2}{1.25}{0.58}}
\ArqMesa{22.5}{2.4}{2.8}{1.35}\ArqBancada{22.5}{6.6}{2.8}{0.45}
% nucleo vertical e sanitarios
\ArqEscada{10.25}{0.35}{2.75}{7.25}
\ArqElevador{13.65}{6.05}{1.75}{1.70}
\draw[arqEquip] (15.60,6.05) rectangle (16.55,7.75);\node[arqSub,align=center] at (16.08,6.90){SHAFT\\TÉCN.};
\ArqVaso{15.45}{0.65}\ArqPia{14.55}{1.18}{0.60}{0.35}\node[arqSub] at (15.45,1.55){WC-M};
\ArqVaso{15.45}{2.55}\ArqPia{14.55}{3.05}{0.60}{0.35}\node[arqSub] at (15.45,3.42){WC-F};
\ArqVaso{15.50}{4.55}\ArqPia{14.55}{5.15}{0.70}{0.38}\node[arqSub] at (15.45,5.55){WC-PCD};
\node[arqSub,rotate=90] at (17.35,3.0){ACESSO WC};
\node[arqSub] at (14.45,7.55){HALL ELEV.};
% rotulos
\foreach \x/\n in {2.6/01,7.6/02,13.0/03,18.2/04,23.4/05}{\node[arqLabel] at (\x,17.35){SALA \n};}
\node[arqLabel] at (2.5,0.65){SALA 06};\node[arqLabel] at (7.5,0.65){SALA 07};\node[arqLabel] at (20.0,0.65){SALA 08};\node[arqLabel] at (24.0,0.65){REUNIÃO};
\node[arqSub] at (5.0,9.0){CORREDOR - 2,00 m};\node[arqSub] at (20.5,9.0){CORREDOR};
\CotaH{0}{26}{18.22}{26,00 m}\CotaV{0}{18}{26.22}{18,00 m}
}

\begin{frame}{Comercial I — pavimento típico refinado}
\centering\begin{tikzpicture}[scale=0.37,every node/.style={font=\tiny}]\EdifComBase\end{tikzpicture}
\vspace{0.4mm}\scriptsize Salas voltadas ao corredor central; escada, elevador e shaft técnico com acessos próprios; sanitários atendidos por corredor lateral de 1,30 m. Potência e telecom usam compartimentação/segregação compatível no shaft.
\end{frame}

\begin{frame}{Comercial I — organização dos quatro pavimentos}
\small\begin{columns}[T,onlytextwidth]
\column{0.24\textwidth}\begin{caixaProjeto}[title=Térreo]\begin{itemize}\item lobby;\item 6 salas;\item reunião;\item WC M/F/PCD;\item elevador + escada.\end{itemize}\end{caixaProjeto}
\column{0.24\textwidth}\begin{caixaProjeto}[title=1º pav.]\begin{itemize}\item 8 salas;\item reunião;\item WC M/F/PCD;\item shaft/IT.\end{itemize}\end{caixaProjeto}
\column{0.24\textwidth}\begin{caixaProjeto}[title=2º pav.]\begin{itemize}\item 8 salas;\item reunião;\item WC M/F/PCD;\item shaft/IT.\end{itemize}\end{caixaProjeto}
\column{0.24\textwidth}\begin{caixaProjeto}[title=3º pav.]\begin{itemize}\item 8 salas;\item reunião;\item WC M/F/PCD;\item acesso técnico.\end{itemize}\end{caixaProjeto}
\end{columns}\end{frame}

\begin{frame}{Comercial I — elétrica do pavimento típico}
\centering\begin{tikzpicture}[scale=0.37,every node/.style={font=\tiny}]
\EdifComBase\SimQD{16.05}{9.65}{QDP}
\foreach \x in {1.3,3.9,6.5,9.1,11.7,14.3,16.9,19.5,22.1,24.7}{\SimLuz{\x}{14.0}}
\foreach \x in {1.3,3.9,6.5,9.1,18.7,20.7,22.7,24.7}{\SimLuz{\x}{4.4}}
\foreach \x in {1.5,4.5,7.5,10.5,13.5,16.5,19.5,22.5,25.0}{\SimLuz{\x}{9.0}}
\foreach \x/\y in {0.4/11,4.8/11,5.6/11,9.8/11,10.8/11,15/11,16/11,20.2/11,21.2/11,25.6/11,0.4/7.2,4.6/7.2,5.4/7.2,9.6/7.2,18.4/7.2,21.5/7.2,22.3/7.2,25.6/7.2}{\SimTUG{\x}{\y}\SimDados{\x}{\y+0.35}}
\foreach \x in {2.6,7.6,13,18.2,23.4}{\SimTUE{\x}{17.6}{AC}}\foreach \x in {2.5,7.5,20,24}{\SimTUE{\x}{0.4}{AC}}\SimTUE{14.5}{7.5}{EL}
\foreach \x in {2,7,12,17,22}{\SimEmerg{\x}{9.0}}\SimEmerg{11.6}{7.55}\SimEmerg{14.5}{7.55}
\draw[corPrim!70,dashed,thick] (16.05,9.65)--(16.05,9)--(2,9);\draw[corPrim!70,dashed,thick] (16.05,9)--(25,9);
\end{tikzpicture}\end{frame}

\begin{frame}[shrink=5]{Comercial I — carga e demanda do estudo}
\scriptsize
\begin{tabularx}{\textwidth}{>{\bfseries}p{3.0cm}c c c c X}
\toprule
Grupo & Dado instalado & $fp$ & $P_{inst}$ & $f_d$ & $P_{dem}$\\
\midrule
Iluminação/emergência & 18,0 kW & 0,95 & 18,0 kW & 0,90 & 16,2 kW\\
TUG/TI & 32,0 kVA & 0,90 & 28,8 kW & 0,70 & 20,2 kW\\
Climatização & 72,0 kW & 0,92 & 72,0 kW & 0,80 & 57,6 kW\\
Elevador & 18,0 kW & 0,85 & 18,0 kW & 0,80 & 14,4 kW\\
Serviços comuns & 20,0 kW & 0,90 & 20,0 kW & 0,60 & 12,0 kW\\
\midrule
\textbf{Total} & -- & -- & \textbf{156,8 kW} & -- & \textbf{120,4 kW}\\
\bottomrule
\end{tabularx}
\begin{caixaAt}
Os fatores de demanda e de potência são \textbf{premissas didáticas}, não uma tabela universal. No executivo, usar ocupação, simultaneidade, dados de placa, estratégia do elevador/HVAC e método aceito pela Equatorial.
\end{caixaAt}
\end{frame}

\begin{frame}{Comercial I — distribuição vertical e conexão}
\small\begin{columns}[T,onlytextwidth]
\column{0.49\textwidth}
\begin{caixaProjeto}[title=Ordem de grandeza]
Pelas premissas do estudo:
\[
S_{dem}\approx\SI{132}{kVA},\qquad
I_{dem}\approx\frac{132000}{\sqrt3\,380}=\SI{201}{A}.
\]
QDI, QDT/TI, QDAC, QD-ELEV e QD-SERV partem do QGBT com alimentadores, proteção e seletividade próprios.
\end{caixaProjeto}
\column{0.49\textwidth}
\begin{caixaNorma}[title=Interface correta]
O empreendimento não pode ser enquadrado pela tabela de unidade individual da NT.00001. Medição coletiva, demanda do conjunto e prumadas seguem a \textbf{NT.00004.EQTL Rev.08}; se a solução de conexão for em MT, aplicar também a \textbf{NT.00002.EQTL Rev.10} e os dados fornecidos pela distribuidora.
\end{caixaNorma}
\end{columns}\end{frame}

\begin{frame}{Comercial I — cobertura e SPDA}
\centering\begin{tikzpicture}[scale=0.36,every node/.style={font=\tiny}]
\draw[arqExt] (0,0) rectangle (26,18);\draw[arqEquip,line width=0.8pt] (10.2,5.2) rectangle (17.2,10.8);\node[arqLabel] at (13.7,8){CASA DE MÁQUINAS\\ACESSO À ESCADA};
\draw[arqEquip] (19,12) rectangle (23.5,15);\node[arqSub] at (21.25,13.5){CONDENSADORAS};
\foreach \y in {0,6,12,18}{\draw[corAt!80,thick] (0,\y)--(26,\y);}\foreach \x in {0,6.5,13,19.5,26}{\draw[corAt!80,thick] (\x,0)--(\x,18);}
\foreach \x/\y in {0/0,6.5/0,13/0,19.5/0,26/0,26/6,26/12,26/18,19.5/18,13/18,6.5/18,0/18,0/12,0/6}{\fill[corNorma] (\x,\y) circle (0.14);}
\end{tikzpicture}\vspace{0.3mm}\scriptsize Malha e descidas são conceituais; classe e espaçamentos finais decorrem da análise de risco e da série ABNT NBR 5419 aplicável.
\end{frame}

\begin{frame}{Comercial I — SPDA, DPS e aterramento integrados}
\centering\begin{tikzpicture}[scale=0.70,every node/.style={font=\tiny}]
\tikzset{b/.style={draw=corPrim,fill=corDef!7,rounded corners,minimum width=2.2cm,minimum height=0.58cm,align=center}}
\node[b](r)at(-4,3.5){Rede Equatorial};\node[b](qg)at(-4,2.3){QGBT\\DPS origem};\node[b](qd)at(-4,1.1){QDs pavimentos};\node[b](sp)at(1,3.5){SPDA cobertura};\node[b](desc)at(1,2.3){Descidas};\node[b](bep)at(1,1.1){BEP / equipot.};\node[b](malha)at(1,-0.1){aterramento};
\draw[-{Stealth},thick,corPrim](r)--(qg)--(qd);\draw[-{Stealth},thick,corAt](sp)--(desc)--(bep)--(malha);\draw[corNorma,thick](qg)--(bep);\draw[corNorma,thick](qd)--(bep);
\end{tikzpicture}\end{frame}

% Comercial II - centro empresarial e de treinamento com dez salas
\section{Projeto Comercial II — centro empresarial e de treinamento}
\newcommand{\CentroServicosBase}{%
\draw[arqExt] (0,0) rectangle (32,20);
\ArqParedeInt{28}{0}{28}{20}\ArqParedeInt{27}{6}{27}{16}
\ArqParedeInt{0}{6}{27}{6}\ArqParedeInt{0}{8}{27}{8}\ArqParedeInt{0}{14}{27}{14}\ArqParedeInt{0}{16}{27}{16}
\ArqParedeInt{13}{0}{13}{6}\ArqParedeInt{19}{0}{19}{6}\ArqParedeInt{23}{0}{23}{6}
% sanitarios lado a lado e lobby proprio
\ArqParedeInt{15}{0}{15}{4.8}\ArqParedeInt{17}{0}{17}{4.8}\ArqParedeInt{13}{4.8}{19}{4.8}
\foreach \x in {4.5,9,13.5,18,22.5}{\ArqParedeInt{\x}{8}{\x}{14}}
\foreach \x in {4,8,12,16,22}{\ArqParedeInt{\x}{16}{\x}{20}}
\ArqParedeInt{28}{7}{32}{7}\ArqParedeInt{28}{12}{32}{12}\ArqParedeInt{30.1}{7}{30.1}{12}
% acessos externos e internos
\ArqPortaCorrerH{4.8}{7.6}{0}\ArqVaoH{10.4}{12.2}{6}\ArqVaoH{15.5}{16.5}{6}
\ArqPortaHDown{13.35}{4.8}{0.80}\ArqPortaHDown{15.35}{4.8}{0.80}\ArqPortaHDown{17.25}{4.8}{0.95}
\ArqPortaHDown{20}{6}{0.90}\ArqPortaHDown{24.3}{6}{0.90}
\foreach \x in {0.7,5.2,9.7,14.2,18.7,23.2}{\ArqPortaHUp{\x}{8}{0.90}}
\foreach \x in {0.6,4.6,8.6,12.6,17.1,23.1}{\ArqPortaHUp{\x}{16}{0.90}}
% corredor vertical liga os dois corredores ao nucleo; escada e elevador abrem para ele
\ArqVaoV{6.35}{7.65}{27}\ArqVaoV{14.35}{15.65}{27}\ArqPortaVRight{28}{6.00}{0.90}\ArqPortaElevV{8.05}{9.05}{28}
% frente
\ArqDesk{9.4}{1.1}{2.5}{0.65}\ArqDesk{9.4}{2.3}{2.5}{0.65}\foreach \x in {1,2.4,3.8,5.2,6.6}{\draw[arqMob] (\x,2) circle (0.28);\draw[arqMob] (\x,3.5) circle (0.28);}\node[arqLabel] at (6.5,5.45){RECEPÇÃO / ESPERA};
\ArqVaso{14}{1.2}\ArqPia{13.25}{3.5}{0.60}{0.38}\node[arqSub] at (14,4.45){WC-M};
\ArqVaso{16}{1.2}\ArqPia{15.25}{3.5}{0.60}{0.38}\node[arqSub] at (16,4.45){WC-F};
\ArqVaso{18}{1.3}\ArqPia{17.25}{3.5}{0.70}{0.40}\node[arqSub] at (18,4.45){WC-PCD};\node[arqSub] at (16,5.45){LOBBY WC};
\ArqDesk{19.6}{1.2}{1.35}{0.58}\node[arqLabel] at (21,5.45){ADMIN.};\ArqBancada{23.5}{0.5}{3.2}{0.55}\ArqMesa{24}{2.5}{2.2}{0.9}\node[arqLabel] at (25,5.45){COPA / APOIO};
% salas 01--10: escritorios/salas de atendimento, sem mobiliario de uso medico
\foreach \x/\n in {0/01,4.5/02,9/03,13.5/04,18/05,22.5/06}{\ArqDesk{\x+0.35}{8.65}{1.35}{0.55}\ArqDesk{\x+2.35}{10.35}{1.35}{0.55}\node[arqLabel] at (\x+2.25,13.55){SALA \n};}
\foreach \x/\n in {0/07,4/08,8/09,12/10}{\ArqDesk{\x+0.25}{16.45}{1.2}{0.5}\ArqDesk{\x+2.15}{18.0}{1.2}{0.5}\node[arqLabel] at (\x+2,19.55){SALA \n};}
\ArqDesk{17.0}{17.0}{1.5}{0.55}\ArqDesk{19.2}{17.0}{1.5}{0.55}\node[arqLabel] at (19,19.55){SALA DE PROJETOS};
\ArqMesa{22.7}{17.0}{3.4}{1.1}\node[arqLabel] at (24.5,19.55){TREINAMENTO / INOVAÇÃO};
% nucleo vertical
\ArqEscada{28.25}{0.35}{3.45}{6.25}\ArqElevador{28.25}{7.35}{1.60}{2.40}\draw[arqEquip] (30.35,7.35) rectangle (31.75,9.75);\node[arqSub] at (31.05,8.55){SHAFT};\node[arqSub,rotate=90] at (27.5,11){CIRC. VERTICAL};
\CotaH{0}{32}{20.22}{32,00 m}\CotaV{0}{20}{32.22}{20,00 m}}

\begin{frame}{Comercial II — centro empresarial com circulação clara}
\centering\begin{tikzpicture}[scale=0.30,every node/.style={font=\tiny}]\CentroServicosBase\end{tikzpicture}
\vspace{0.4mm}\scriptsize Salas de trabalho, recepção, apoio e treinamento têm acessos próprios. O lobby dos sanitários é separado, e a circulação vertical conecta diretamente os corredores à escada e ao elevador. O caso não é tratado como estabelecimento assistencial de saúde.
\end{frame}

\begin{frame}{Comercial II — iluminação normal e de emergência}
\centering\begin{tikzpicture}[scale=0.30,every node/.style={font=\tiny}]
\CentroServicosBase\SimQD{27.25}{7.20}{QDC}
\foreach \x in {2,5,8,11,14,17,20,23,26}{\SimLuz{\x}{7}\SimLuz{\x}{15}}
\foreach \x/\y in {3/3,7/3,11/3,2.25/11,6.75/11,11.25/11,15.75/11,20.25/11,24.75/11,2/18,6/18,10/18,14/18,19/18,24.5/18}{\SimLuz{\x}{\y}}
\foreach \x in {2,7,12,17,22,26}{\SimEmerg{\x}{7.0}\SimEmerg{\x}{15.0}}\SimEmerg{29.8}{3.0}\SimEmerg{29.0}{8.5}
\draw[corAccent!75!black,dashed,thick] (27.25,7.20)--(2,7);\draw[corAccent!75!black,dashed,thick] (27.25,7.20)--(27.25,15)--(2,15);
\end{tikzpicture}
\vspace{0.3mm}\scriptsize Circuitos independentes para iluminação normal, emergência/rotas de fuga e áreas comuns; autonomia e níveis finais dependem do projeto de segurança contra incêndio aplicável.
\end{frame}

\begin{frame}{Comercial II — TUG, dados, climatização e serviços}
\centering\begin{tikzpicture}[scale=0.30,every node/.style={font=\tiny}]
\CentroServicosBase\SimQD{27.25}{7.20}{QDC}
% duas estações TUG+dados por sala 01--06
\foreach \x in {0.55,5.05,9.55,14.05,18.55,23.05}{\SimTUG{\x}{9.0}\SimDados{\x}{9.45}\SimTUG{\x+3.35}{12.8}\SimDados{\x+3.35}{12.35}}
% salas 07--10, projetos e treinamento
\foreach \x in {0.45,4.45,8.45,12.45}{\SimTUG{\x}{16.6}\SimDados{\x}{17.05}\SimTUG{\x+3.1}{19.2}\SimDados{\x+3.1}{18.75}}
\foreach \x/\y in {17.2/16.6,20.8/19.2,22.8/16.6,26.4/19.2,1.0/1.0,6.5/4.8,10.0/1.0,20.0/1.0}{\SimTUG{\x}{\y}\SimDados{\x}{\y+0.45}}
\foreach \x/\y in {23.7/0.45,24.8/0.45,26.0/0.45,26.6/3.5}{\SimTUGM{\x}{\y}}
\foreach \x in {2.25,6.75,11.25,15.75,20.25,24.75}{\SimTUE{\x}{13.7}{AC}}
\foreach \x in {2,6,10,14,19,24.5}{\SimTUE{\x}{19.7}{AC}}
\SimTUE{8}{0.45}{AC}\SimTUE{21}{0.45}{AC}\SimTUE{29.0}{8.5}{EL}
\draw[corPrim!70,dashed,thick] (27.25,7.20)--(2,7);\draw[corPrim!70,dashed,thick] (27.25,7.20)--(27.25,15)--(2,15);
\end{tikzpicture}
\vspace{0.2mm}\tiny Os símbolos representam conjuntos funcionais; quantidade de tomadas e portas de dados deve fechar com postos de trabalho, salas de treinamento, copa e reserva de expansão.
\end{frame}

\begin{frame}[shrink=5]{Comercial II — carga e demanda do estudo}
\scriptsize
\begin{tabularx}{\textwidth}{>{\bfseries}p{3.0cm}c c c c X}
\toprule
Grupo & Dado instalado & $fp$ & $P_{inst}$ & $f_d$ & $P_{dem}$\\
\midrule
Iluminação/emergência & 14,0 kW & 0,95 & 14,0 kW & 0,90 & 12,6 kW\\
TUG/dados/TI & 36,0 kVA & 0,90 & 32,4 kW & 0,65 & 21,1 kW\\
Climatização & 64,0 kW & 0,92 & 64,0 kW & 0,80 & 51,2 kW\\
Elevador & 18,0 kW & 0,85 & 18,0 kW & 0,80 & 14,4 kW\\
Copa e serviços & 16,0 kW & 0,90 & 16,0 kW & 0,60 & 9,6 kW\\
\midrule
\textbf{Total} & -- & -- & \textbf{144,4 kW} & -- & \textbf{108,9 kW}\\
\bottomrule
\end{tabularx}
\begin{caixaAt}
Valores e fatores são hipóteses didáticas rastreáveis. O executivo substitui todos eles por equipamentos, ocupação, simultaneidade e método aceito pela distribuidora.
\end{caixaAt}
\end{frame}

\begin{frame}{Comercial II — distribuição e enquadramento correto}
\small\begin{columns}[T,onlytextwidth]
\column{0.49\textwidth}
\begin{caixaProjeto}[title=Distribuição]
\[
S_{dem}\approx\SI{120}{kVA},\qquad I_{dem}\approx\SI{182}{A}\ @\ \SI{380}{V}.
\]
Separar QDL, QDT/TI, QDAC, QD-ELEV e QD-SERV; verificar queda, curto, seletividade, reserva e continuidade requerida.
\end{caixaProjeto}
\column{0.49\textwidth}
\begin{caixaNorma}[title=Conexão]
Como a carga instalada supera 75 kW, não usar a tabela de unidade individual da NT.00001. Se houver uma única UC, solicitar estudo e aplicar a NT.00002 quando o atendimento for em MT; se o prédio for subdividido em várias UCs, aplicar a NT.00004 ao EMUC.
\end{caixaNorma}
\end{columns}
\end{frame}

\begin{frame}{Comercial II — unifilar funcional}
\centering\begin{tikzpicture}[scale=0.90,every node/.style={font=\scriptsize}]
\tikzset{b/.style={draw=corPrim,fill=corDef!8,rounded corners,minimum width=2.15cm,minimum height=0.58cm,align=center}}
\node[b](orig)at(0,4.2){Origem definida\\pela Equatorial};\node[b](qg)at(0,3.2){QGBT\\DG + DPS};
\node[b](l)at(-4,1.4){QDL\\normal/emergência};\node[b](t)at(-2,0.2){QDT/TI\\TUG + dados};\node[b](ac)at(0,1.4){QDAC\\climatização};\node[b](el)at(2,0.2){QD-ELEV\\elevador};\node[b](sv)at(4,1.4){QD-SERV\\copa/serviços};
\node[b,draw=corNorma,fill=corNorma!7](bep)at(4.4,3.2){BEP\\PE/equipot.};
\draw[-{Stealth},thick,corPrim](orig)--(qg);\foreach \x in {l,t,ac,el,sv}{\draw[-{Stealth},thick](qg)--(\x);}\draw[corNorma,thick](qg)--(bep);\draw[corNorma,thick](bep.south)--(4.4,2.45)node[ground]{};
\end{tikzpicture}
\end{frame}

\input{projetos/industriais/industria_bt.tex}
\input{projetos/industriais/industria_subestacao.tex}
\section{Comparação e entrega executiva}

\begin{frame}{Projetos do caderno — residenciais e comerciais}
\scriptsize
\begin{tabularx}{\textwidth}{c p{3.0cm} p{3.1cm} X}
\toprule
ID & Projeto & Diferencial de arquitetura & Foco elétrico \\
\midrule
R1 & Casa térrea & 2 suítes, banho social, quarto 3, garagem e hall interno & entrada BT, balanceamento e três chuveiros \\
R2 & Apartamento & 2 suítes, varanda, lavabo e serviço & unidade em EMUC, prumada e QD trifásico \\
C1 & Edifício corporativo & 4 pavimentos, 8 salas/tipo, WC M/F/PCD, elevador, escada & QGBT, elevador, SPDA, DPS \\
C2 & Centro empresarial & recepção, 10 salas, projetos, treinamento, copa e núcleo vertical & QGBT, TI, climatização, elevador e emergência \\
C3 & Loja de shopping & área de vendas, vitrines, provadores, caixa, estoque e sala elétrica & luminotécnica, teto refletido, cenas, alimentador e comissionamento \\
\bottomrule
\end{tabularx}
\end{frame}

\begin{frame}{Projetos do caderno — industriais}
\scriptsize
\begin{tabularx}{\textwidth}{c p{3.0cm} p{3.1cm} X}
\toprule
ID & Projeto & Diferencial de arquitetura & Foco elétrico \\
\midrule
I1 & Indústria leve & produção A/B, manutenção, almoxarifado, vestiários & motores/VFD/CCM com atendimento em BT \\
I2 & Indústria com SE & processo, embalagem, expedição, utilidades e refrigeração & MT, trafo 500 kVA, QGBT 1.000 A e SPDA \\
\bottomrule
\end{tabularx}
\begin{caixaProjeto}[title={Sete casos, quatro competências}]
Em todos os projetos, o aluno deve localizar e interpretar informações, dimensionar a solução, compreender a ligação/montagem e planejar execução, ensaios e entrega.
\end{caixaProjeto}
\end{frame}

\begin{frame}{O que muda de um projeto para outro}
\small
\begin{columns}[T,onlytextwidth]
\column{0.49\textwidth}
\begin{caixaProjeto}[title=Não são a mesma planta reciclada]
\begin{itemize}
\item geometrias e circulações diferentes;
\item usos e cargas próprios;
\item critérios luminotécnicos e cenas próprios na loja;
\item quadros em posições coerentes;
\item rotas compatíveis com cada arquitetura;
\item distribuição e unifilares diferentes.
\end{itemize}
\end{caixaProjeto}
\column{0.49\textwidth}
\begin{caixaNorma}[title=Norma conforme o caso]
\begin{itemize}
\item BT individual: NT.00001.EQTL;
\item múltiplas unidades: NT.00004.EQTL;
\item MT/subestação: NT.00002 e NT.00026;
\item SPDA: série ABNT NBR 5419;
\item loja: ABNT NBR ISO/CIE 8995-1, NBR 10898 e projeto de incêndio aprovado;
\item MT interna: ABNT NBR 14039.
\end{itemize}
\end{caixaNorma}
\end{columns}
\end{frame}

\begin{frame}{Checklist de projeto executivo}
\small
\begin{enumerate}
\item Usar a planta arquitetônica real, com cotas, níveis, layout e equipamentos.
\item Calcular iluminação/TUG e levantar TUE por dados de placa.
\item Dividir circuitos e definir fases.
\item Calcular $I_b$, escolher seção por $I_z$ e fatores de correção.
\item Verificar $I_b\le I_n\le I_z$ e condições de proteção.
\item Calcular queda de tensão, curto máximo/mínimo e seletividade.
\item Dimensionar N, PE, eletrodutos, bandejas, quadros e barramentos.
\item Integrar DR, DPS, aterramento e SPDA sem criar aterramentos independentes artificiais.
\item Em MT, solicitar dados de rede e coordenar proteção da Equatorial com a instalação interna.
\item Emitir memorial, documentação de responsabilidade técnica, comissionamento e as built.
\end{enumerate}
\end{frame}

\section{Referências e rastreabilidade}

\begin{frame}[shrink=4]{Fontes oficiais verificadas nesta revisão}
\small
\begin{itemize}
\item \textbf{ABNT Catálogo}: confirmação da ABNT NBR 5419-1:2026 e consulta das edições indicadas.\\
\href{https://www.abntcatalogo.com.br/}{\texttt{abntcatalogo.com.br}}
\item \textbf{Equatorial Energia Maranhão}: lista vigente das normas de conexão e respectivos anexos, incluindo NT.00001 Rev.09, NT.00002 Rev.10, NT.00004 Rev.08, NT.00020 Rev.06, NT.00030 Rev.03 e NT.00042 Rev.04.\\
\href{https://ma.equatorialenergia.com.br/institucional/leis-e-normas-tecnicas/normas-tecnicas/instalacao-de-padrao/}{\texttt{Equatorial MA — Normas Técnicas de Conexão}}
\item \textbf{Ministério do Trabalho e Emprego}: NR-10 vigente até 31/05/2027 e nova redação com vigência em 01/06/2027.\\
\href{https://www.gov.br/trabalho-e-emprego/pt-br/acesso-a-informacao/participacao-social/conselhos-e-orgaos-colegiados/comissao-tripartite-partitaria-permanente/normas-regulamentadora/normas-regulamentadoras-vigentes/norma-regulamentadora-no-10-nr-10}{\texttt{MTE — página oficial da NR-10}}
\end{itemize}
\begin{caixaAt}
Consulta realizada em 13/08/2026. A lista da distribuidora pode mudar; no início de cada projeto, conferir novamente revisão, anexos e período de transição. Este material resume critérios e não reproduz nem substitui o texto integral das normas ABNT.
\end{caixaAt}
\end{frame}

\begin{frame}{Referências para a loja de shopping}
\small
\begin{itemize}
\item \textbf{ABNT NBR ISO/CIE 8995-1}: iluminação de ambientes de trabalho interiores; confirmar edição vigente no \href{https://www.abntcatalogo.com.br/}{Catálogo ABNT}.
\item \textbf{ABNT NBR 10898}: sistema de iluminação de emergência; usar em conjunto com o projeto de segurança contra incêndio aprovado.
\item \textbf{CBMMA}: \href{https://cbm.ssp.ma.gov.br/unidades-bm/unidades-administrativas/dat/legislacao-tecnica-dat/}{legislação técnica oficial}, incluindo NT 18 para iluminação de emergência e demais normas aplicáveis ao shopping.
\item \textbf{Manual técnico do empreendimento}: ponto de entrega interno, demanda concedida, proteção, medição, materiais, procedimentos de obra e comissionamento.
\item \textbf{Dados fotométricos e fichas dos fabricantes}: arquivos IES/LDT, fluxo, distribuição, potência, IRC, CCT, compatibilidade de controles, torque e manutenção.
\end{itemize}
\begin{caixaAt}
As metas numéricas do estudo da loja são premissas didáticas declaradas. O executivo deve confrontá-las com as edições normativas vigentes, a atividade visual real, a marca e as exigências do shopping.
\end{caixaAt}
\end{frame}


\begin{frame}[plain]
\begin{tikzpicture}[remember picture,overlay]
\fill[corPrim] (current page.south west) rectangle (current page.north east);
\fill[corAccent] (current page.south west) rectangle ++(0.18\paperwidth,\paperheight);
\node[anchor=west,white,font=\fontsize{31}{35}\bfseries,text width=11cm] at ([xshift=2.5cm,yshift=0.6cm]current page.west) {Fim do caderno de projetos};
\node[anchor=west,corAccent!30!white,font=\fontsize{14}{17}\bfseries,text width=11.5cm,align=left] at ([xshift=2.55cm,yshift=-0.5cm]current page.west) {Seis projetos distintos — Equatorial Maranhão};
\end{tikzpicture}
\end{frame}
\end{document}
